\chapter{Literature Survey}
\label{ch:2}


\begin{figure}[H]
\centering
\includegraphics[width=0.8\textwidth]{10_federated_vs_centralized.png}
\caption{Federated vs Centralized Learning Comparison}
\label{fig:10_federated_vs_centralized}
\end{figure}


\section{2.1 Introduction to Literature Context}
Developing a privacy-preserving brain tumor classification system requires understanding three interconnected research domains:

\begin{enumerate}
  \item **Federated Learning (FL):** Algorithms, optimization strategies, and non-IID challenges
  \item **Medical Image Analysis:** Brain tumor classification, segmentation, and progression forecasting
  \item **Privacy and Security in Healthcare:** Regulatory frameworks and technical implementations
\end{enumerate}

This chapter synthesizes foundational work, identifies research gaps, and positions this project's contributions within the broader context.


\section{2.2 Federated Learning: Fundamentals and Evolution}

\subsection{2.2.1 Foundational Algorithm: Federated Averaging (FedAvg)}
McMahan et al. (2017) introduced Federated Averaging, the foundational algorithm that revolutionized distributed learning. In their seminal paper, they demonstrated that training a global model across multiple decentralized devices (clients) through weighted averaging could achieve comparable accuracy to centralized training while preserving privacy.

Key Contributions:

\begin{itemize}
  \item **Algorithm:** Simple weighted averaging of client-side model weights by dataset size
  \item **Convergence Proof:** Theoretical convergence guarantees under convex loss functions
  \item **Practical Impact:** Enabled large-scale distributed training on mobile devices, IoT sensors, and healthcare institutions
\end{itemize}

Formula:

w\_global\^(t+1) = Σ (n\_k / N) × w\_k\^(t)

where n\_k is Client k's dataset size, N is total data across all clients.

Limitations Identified: FedAvg performs poorly when client data is non-IID (heterogeneous class distributions, different data qualities across institutions).


\subsection{2.2.2 Non-IID Challenge and Proximal Solutions}
Li et al. (2020) extended FedAvg by introducing Federated Proximal (FedProx), addressing the core challenge of data heterogeneity in federated settings.

Problem: When clients have different data distributions (non-IID), naively averaging weights from diverse clients can result in:

\begin{itemize}
  \item Rapid divergence of client models from the global optimum
  \item Slower convergence
  \item Bias toward clients with larger datasets
\end{itemize}

FedProx Solution:

\begin{itemize}
  \item Added a **proximal regularization term** to local client training loss
  \item Prevents clients from drifting too far from the global model during local updates
  \item Mathematically: `loss\_local = loss\_data + (μ/2) × ||w\_local - w\_global||²`
\end{itemize}

Clinical Relevance: In multi-hospital settings where each hospital has different patient populations (e.g., different tumor incidence rates, imaging protocols), FedProx helps maintain model agreement while allowing local adaptation.

Limitation: Proximal term can be too restrictive, preventing clients from learning local-specific patterns (explored in this project through fairness analysis).


\subsection{2.2.3 Variance Reduction: SCAFFOLD}
Karimireddy et al. (2020) introduced SCAFFOLD (Stochastic Controlled Averaging for Federated Learning), addressing client-drift through variance reduction:

\begin{itemize}
  \item Maintains client control variates to track divergence from global model
  \item Reduces impact of heterogeneous data on convergence
  \item Provides improved convergence guarantees for non-convex functions (relevant for neural networks)
\end{itemize}

Not used in this project but important context: SCAFFOLD demonstrates that sophisticated variance-based methods can improve FL robustness.


\subsection{2.2.4 Non-IID Analysis and Characterization}
Zhao et al. (2018) conducted foundational analysis of non-IID data in federated settings, categorizing heterogeneity types:

This Project: Implements both natural heterogeneity (Setup 1) and moderate label skew (Setup 2) based on Zhao et al.'s taxonomy.


\section{2.3 Quantum-Inspired Optimization: QPSO}

\subsection{2.3.1 Classical Particle Swarm Optimization (PSO)}
Kennedy \& Eberhart (1995) introduced Particle Swarm Optimization, a bio-inspired optimization algorithm mimicking bird flocking behavior. Each "particle" represents a potential solution, iteratively improving through:

\begin{itemize}
  \item Personal best memory (pbest): best position each particle has found
  \item Global best memory (gbest): best position found by any particle
  \item Velocity update rules balancing exploration and exploitation
\end{itemize}

PSO Formula:

v(t+1) = w×v(t) + c1×r1×(pbest - x(t)) + c2×r2×(gbest - x(t))
x(t+1) = x(t) + v(t+1)


\subsection{2.3.2 Quantum-Behaved PSO (QPSO)}
Sun et al. (2004, 2012) introduced Quantum-Behaved Particle Swarm Optimization, replacing classical velocity dynamics with quantum mechanics-inspired position updates:

Key Innovation: Instead of classical velocity, particles occupy quantum states with probabilistic position distributions. Position update:

x\_new = p ± β × |mbest - x| × ln(1/u)

where:

\begin{itemize}
  \item `p` is the attraction point (weighted combination of pbest and gbest)
  \item `β` is the contraction-expansion coefficient (exploration vs. exploitation)
  \item `u ∈ [0,1]` is a random variable providing stochasticity
  \item `mbest` is the mean of all personal bests
\end{itemize}

Advantages Over Classical PSO:

\begin{itemize}
  \item **Faster convergence:** Quantum uncertainty principle enables more efficient exploration
  \item **Better escape from local optima:** Stochastic quantum jumps explore solution space more thoroughly
  \item **Fewer parameters:** Only `β` and `u` range need tuning vs. 3+ parameters in classical PSO
\end{itemize}

Application to Federated Learning (Novel in this Project):

\begin{itemize}
  \item Treat each client's model weights as a "particle" in optimization space
  \item Track personal best (best weights client has produced) and global best (best weights overall)
  \item Use QPSO updates to aggregate weights, balancing stability (exploitation) with exploration of better solutions
\end{itemize}


\section{2.4 Brain Tumor Classification and Detection}

\subsection{2.4.1 Traditional CNN Approaches}
Krizhevsky et al. (2012) introduced AlexNet, demonstrating deep convolutional neural networks' superiority on image classification tasks. For medical imaging:

ResNet Architecture (He et al., 2015):

\begin{itemize}
  \item Introduced residual connections enabling training of very deep networks
  \item ResNet-18: 18 layers, \~11.2M parameters
  \item ImageNet pretraining provides strong feature initialization for downstream tasks like brain tumor classification
\end{itemize}

In this Project: ResNet-18 serves as the classification backbone across all 3 federated clients, fine-tuned for 3-class (Glioma, Meningioma, Pituitary) and 4-class problems.


\subsection{2.4.2 Transfer Learning for Medical Imaging}
Tan \& Le (2019) systematized transfer learning, showing pre-trained features from large-scale datasets (ImageNet) transfer effectively to specialized medical domains with limited data.

Findings Relevant to Our Work:

\begin{itemize}
  \item Pre-trained ImageNet weights significantly accelerate medical image model training
  \item Fine-tuning on task-specific data improves performance
  \item Works well even with moderate dataset sizes (9,000+ images as in this project)
\end{itemize}


\subsection{2.4.3 Brain Tumor Classification Datasets}
Masoud Brain Tumor MRI (Kaggle Dataset - nickparvar):

\begin{itemize}
  \item \~7,000 2D MRI slices
  \item Classes: Glioma, Meningioma, No Tumor, Pituitary
  \item Splits available: training/testing
  \item Used in Client 1 and Client 3 in federated setup
\end{itemize}

BRISC 2025 (Brain Tumor Image Segmentation Challenge):

\begin{itemize}
  \item \~4,600 multimodal 3D brain MRI volumes
  \item Standardized preprocessing and evaluation metrics
  \item Used in Client 2 federated setup
\end{itemize}

Masoud + BRISC Combination:

\begin{itemize}
  \item Provides natural data heterogeneity: different imaging protocols, scanner types, patient populations
  \item Simulates realistic multi-hospital scenario without explicit dataset redistribution
\end{itemize}


\section{2.5 Brain Tumor Segmentation}

\subsection{2.5.1 U-Net: Foundational Architecture}
Ronneberger et al. (2015) introduced U-Net, the seminal architecture for biomedical image segmentation:

Architecture Features:

\begin{itemize}
  \item **Encoder-Decoder Structure:** Downsampling path captures context, upsampling path enables precise localization
  \item **Skip Connections:** Concatenate encoder features to decoder, preserving fine-grained spatial information
  \item **Efficiency:** Requires relatively few training samples compared to fully convolutional networks
\end{itemize}

Extension to 3D: MONAI framework implements 3D variants for volumetric medical imaging.


\subsection{2.5.2 Attention Mechanisms in U-Net}
Oktay et al. (2018) enhanced U-Net with Attention Gates (AGs), enabling the network to automatically learn which spatial regions are important:

Attention Gate Mechanism:

\begin{itemize}
  \item Suppresses irrelevant regions (healthy tissue)
  \item Focuses network capacity on target regions (tumor)
  \item Mathematically: `attention\_coeff = sigmoid(transform(encoder\_features, decoder\_features))`
  \item Applied channel-wise, reducing computational overhead
\end{itemize}

In this Project: 3D Attention U-Net processes 4-channel MRI inputs (T1, T1ce, T2, FLAIR) to segment 3 tumor sub-regions (WT, TC, ET).


\subsection{2.5.3 BraTS Challenge and Benchmarks}
Menze et al. (2014) → Baid et al. (2021): Multi-year Brain Tumor Segmentation Challenge establishing benchmarks and best practices:

This Project: Uses BraTS 2021 for Module 2 (Segmentation), achieving Mean Dice 0.76 against benchmark standards (\~0.75-0.78).


\section{2.6 Tumor Progression and Growth Forecasting}

\subsection{2.6.1 Mathematical Growth Models}
Classical mathematical models from population biology have been adapted for tumor growth:

Logistic Growth Model (Verhulst, 1838 → Adapted for Tumors):

V(t) = K / (1 + ((K - V0)/V0) × exp(-r×t))

\begin{itemize}
  \item `K`: Carrying capacity (maximum sustainable tumor volume)
  \item `V0`: Initial tumor volume
  \item `r`: Intrinsic growth rate
  \item Produces S-shaped curve with growth saturation
\end{itemize}

Gompertz Growth Model (Gompertz, 1825 → Adapted for Tumors):

V(t) = V0 × exp((a/b) × (1 - exp(-b×t)))

\begin{itemize}
  \item More realistic for solid tumors (deceleration over time)
  \item Biologically motivated: growth rate decreases as tumor size increases
\end{itemize}

Clinical Application: Different tumor types show different growth patterns:

\begin{itemize}
  \item **HGG (High-Grade Glioma):** Aggressive, exponential-like growth
  \item **LGG (Low-Grade Glioma):** Slow, often logistic or gompertz patterns
\end{itemize}


\subsection{2.6.2 LSTM for Time-Series Medical Data}
Hochreiter \& Schmidhuber (1997): Introduced Long Short-Term Memory networks, capable of learning long-range dependencies in sequential data.

LSTM Architecture:

\begin{itemize}
  \item **Input Gate:** Controls information flow into cell state
  \item **Forget Gate:** Selectively retains or discards past information
  \item **Output Gate:** Controls what cell state information is output
\end{itemize}

Medical Time-Series Applications:

\begin{itemize}
  \item Patient vital sign monitoring
  \item Disease progression prediction
  \item Treatment response forecasting
\end{itemize}


\subsection{2.6.3 Hybrid Approaches: Mathematical + Deep Learning}
Recent Trend in Medical AI (Rajkomar et al., 2018):

\begin{itemize}
  \item Combine interpretable mathematical models with flexible deep learning
  \item Approach: Hybrid = Math\_baseline + LSTM\_correction
  \item Allows model interpretability (math component) + adaptability (LSTM)
\end{itemize}

In this Project:

V\_hybrid(t) = V\_logistic(t) + LSTM\_correction(residual\_sequence)

Achieves 7.88\% MAE improvement for HGG patients by learning where logistic model fails.


\section{2.7 Privacy and Security in Healthcare AI}

\subsection{2.7.1 Regulatory Framework}
HIPAA (Health Insurance Portability and Accountability Act, USA):

\begin{itemize}
  \item Requires de-identification or aggregate statistical analysis for research
  \item Prohibits direct patient data sharing without explicit consent
\end{itemize}

GDPR (General Data Protection Regulation, EU):

\begin{itemize}
  \item Right to deletion ("right to be forgotten")
  \item Data minimization principle
  \item Data processing transparency
\end{itemize}

Challenge: Traditional federated learning still requires hospitals to trust a central server. With GDPR, even aggregated insights may require consent.


\subsection{2.7.2 Differential Privacy in FL}
Dwork et al. (2006): Introduced formal definition of Differential Privacy:

Definition: An algorithm is ε-differentially private if removing one patient's data changes output probability by at most e\^ε.

Application to FL:

\begin{itemize}
  \item Add Laplace or Gaussian noise to client gradients before sending to server
  \item Provides formal privacy guarantees independent of computational power of adversary
  \item Trade-off: privacy vs. model accuracy
\end{itemize}

Not Implemented in This Project: Module 1 implements structural privacy (FL) but not formal DP-epsilon guarantees. Recommended future work.


\subsection{2.7.3 Secure Aggregation}
Bonawitz et al. (2017): Introduced secure aggregation protocol where server cannot decode individual client updates, only aggregated result.

Key Idea:

\begin{itemize}
  \item Client weights encrypted before transmission
  \item Server performs encrypted aggregation (uses homomorphic encryption properties)
  \item Only aggregated result decrypted
\end{itemize}

Implementation Complexity: Trade-off with computational overhead. Not implemented in this project (synchronous aggregation without encryption).


\section{2.8 Research Gaps and Project Positioning}

\subsection{2.8.1 Identified Gaps}

\subsection{2.8.2 Positioning}
This project sits at the intersection of three research areas:

Novelty:

\begin{itemize}
  \item First application of QPSO to federated medical imaging
  \item Fairness-aware aggregation strategy proven statistically superior (p = 2.91 × 10⁻²²) in non-IID settings
  \item Integrated framework allowing multi-hospital brain tumor analysis with privacy guarantees
\end{itemize}


\section{2.9 Chapter Summary}

\section{2.10 Key References}
Literature Cited:

\begin{enumerate}
  \item **McMahan et al. (2017)** - "Communication-Efficient Learning of Deep Networks from Decentralized Data"
\end{enumerate}

AISTATS 2017 — Introduced FedAvg algorithm

\begin{enumerate}
  \item **Li et al. (2020)** - "Federated Optimization in Heterogeneous Networks"
\end{enumerate}

MLSys 2020 — Introduced FedProx for non-IID data

\begin{enumerate}
  \item **Karimireddy et al. (2020)** - "SCAFFOLD: Stochastic Controlled Averaging for Federated Learning"
\end{enumerate}

ICML 2020 — Variance reduction approach

\begin{enumerate}
  \item **Zhao et al. (2018)** - "Federated Learning with Non-IID Data"
\end{enumerate}

arXiv:1806.00582 — Analysis of data heterogeneity

\begin{enumerate}
  \item **Sun et al. (2004, 2012)** - "Quantum-Behaved Particle Swarm Optimization"
\end{enumerate}

IJCNN 2004, IEEE Transactions 2012 — Quantum PSO algorithm

\begin{enumerate}
  \item **He et al. (2015)** - "Deep Residual Learning for Image Recognition"
\end{enumerate}

CVPR 2015 — ResNet architecture

\begin{enumerate}
  \item **Ronneberger et al. (2015)** - "U-Net: Convolutional Networks for Biomedical Image Segmentation"
\end{enumerate}

MICCAI 2015 — U-Net segmentation architecture

\begin{enumerate}
  \item **Oktay et al. (2018)** - "Attention U-Net: Learning Where to Look for the Pancreas"
\end{enumerate}

MIDL 2018 — Attention mechanisms for medical imaging

\begin{enumerate}
  \item **Baid et al. (2021)** - "The RSNA-ASNR-MICCAI BraTS 2021 Benchmark on Brain Tumor Segmentation and Radiogenomic Classification"
\end{enumerate}

arXiv:2107.02314 — BraTS 2021 benchmark dataset

\begin{enumerate}
  \item **Hochreiter \& Schmidhuert (1997)** - "Long Short-Term Memory"
\end{enumerate}

Neural Computation 9(8) — LSTM architecture

\begin{enumerate}
  \item **Dwork et al. (2006)** - "Differential Privacy"
\end{enumerate}

ICALP 2006 — Formal privacy framework

\begin{enumerate}
  \item **Bonawitz et al. (2017)** - "Towards Federated Learning at Scale"
\end{enumerate}

MLSys 2019 — Secure aggregation protocol

\begin{enumerate}
  \item **Edla (2025)** - "Enhancing Federated Learning with Quantum-Inspired PSO: An IID MNIST Study"
\end{enumerate}

Matrusri Engineering College — Prior work on QPSO for FL


\section{Chapter Summary}
This chapter reviewed existing literature on federated learning, brain tumor analysis, and identified key research gaps in fairness-aware FL for healthcare.
